\documentclass[11pt,a4paper]{article}
\usepackage[utf8]{inputenc}
\usepackage[spanish,es-tabla]{babel}
\usepackage[margin=2.2cm]{geometry}
\usepackage{amsmath,amssymb}
\usepackage{booktabs}
\usepackage{caption}
\usepackage{microtype}
\usepackage{tikz}
\usetikzlibrary{shapes.geometric, arrows.meta, positioning, calc}
\usepackage{hyperref}

\hypersetup{
    colorlinks=true,
    linkcolor=black,
    citecolor=black,
    urlcolor=black
}

% Prevent hyphenation in TikZ boxes
\hyphenpenalty=10000
\exhyphenpenalty=10000

\title{\textbf{Proyecto Júpiter: Sistema de Alerta Temprana e Inteligencia Hidrometeorológica para La Serena}\\
\large Informe Técnico de Funcionamiento y Simulación Operativa}
\author{\textbf{Moisés Amundarain}\\
Investigador Principal --- Proyecto Júpiter}
\date{30 de julio de 2026}

\begin{document}
\spanishdeactivate{><}

\maketitle

\begin{abstract}
El presente documento describe la arquitectura, el fundamento hidrológico y el funcionamiento operativo del \textbf{Proyecto Júpiter}, una plataforma predictiva de alerta temprana ante inundaciones y aluviones en la comuna de La Serena. El sistema integra telemetría satelital en tiempo casi real (NRT), un modelo de aprendizaje automático (\textit{Random Forest}) calibrado con física de suelos (método SCS-CN) y un escáner de 20 micro-zonas geográficas calibradas en coordenadas WGS84. Asimismo, se expone la simulación operativa para el evento meteorológico del viernes 31 de julio, detallando el cálculo anticipado de los picos de precipitación y la determinación de las ventanas de calma y retorno seguro.
\end{abstract}

\section{Introducción y Contexto}

El \textbf{Proyecto Júpiter} recibe su nombre en homenaje a \textit{Júpiter Pluvio}, deidad de la mitología romana asociada a las nubes y las lluvias propiciadoras de vida. La plataforma nace para resolver una vulnerabilidad histórica en la comuna de La Serena: la falta de sistemas de alerta con alta resolución espacial y capacidad de anticipación temporal.

Las crecidas repentinas en quebradas precordilleranas (como Pueblo Islón, Lambert y Las Rojas) y los colapsos de colectores pluviales en la zona urbana ocurren debido a la interacción entre tres factores principales:
\begin{enumerate}
    \item Precipitación intensa acumulada en breves periodos.
    \item Límite de nieve (isoterma cero) elevado en la cordillera, que impide la acumulación de nieve y transforma la precipitación en agua líquida sobre la roca.
    \item Saturación previa del suelo, la cual elimina la capacidad de infiltración del terreno.
\end{enumerate}

\section{Arquitectura del Sistema}

El sistema opera bajo un flujo continuo divido en tres capas fundamentales. A diferencia de los pronósticos generales que evalúan a toda la Región de Coquimbo como un único bloque, Proyecto Júpiter divide a La Serena en 20 micro-zonas geográficas independientes con coordenadas WGS84 de alta precisión, tales como Pueblo Islón ($-29.878^\circ, -71.218^\circ$) y Lambert ($-29.818^\circ, -71.148^\circ$).

La estructura de flujo se ilustra en la Figura~\ref{fig:arquitectura}.

\begin{figure}[htbp]
\centering
\begin{tikzpicture}[node distance=1.2cm and 1.5cm, auto, >=Stealth]
    % Fila superior
    \node [draw, rectangle, rounded corners=5pt, text width=3.4cm, minimum height=1.6cm, align=center] (ingesta) 
        {\textbf{1. Ingesta Satelital}\\ Telemetría NRT \\ cada 5 minutos};
    
    \node [draw, rectangle, rounded corners=5pt, text width=3.6cm, minimum height=1.6cm, right=of ingesta, align=center] (features) 
        {\textbf{2. Ingeniería de Datos}\\ Humedad Suelo ($API$)\\ Infiltración (SCS-CN)};
    
    \node [draw, rectangle, rounded corners=5pt, text width=3.6cm, minimum height=1.6cm, right=of features, align=center] (modelo) 
        {\textbf{3. Modelo de IA}\\ Ensamble Classifier \\ Random Forest};
    
    % Fila inferior (sin colisiones de flechas)
    \node [draw, rectangle, rounded corners=5pt, text width=4.2cm, minimum height=1.6cm, below=1.2cm of features, align=center] (scanner) 
        {\textbf{4. Escáner Espacial}\\ 20 Micro-Zonas WGS84 \\ Llegada Pico y Calma};
    
    \node [draw, rectangle, rounded corners=5pt, text width=4.2cm, minimum height=1.6cm, below=1.2cm of modelo, align=center] (dashboard) 
        {\textbf{5. Consola Táctica}\\ Mapa Interactivo \\ Bitácora SCI-201};

    % Conexiones claras y limpias
    \draw [->, line width=1.1pt] (ingesta) -- (features);
    \draw [->, line width=1.1pt] (features) -- (modelo);
    \draw [->, line width=1.1pt] (modelo) -- (dashboard);
    \draw [->, line width=1.1pt] (features) -- (scanner);
    \draw [->, line width=1.1pt] (scanner) -- (dashboard);
\end{tikzpicture}
\caption{Flujo general de procesamiento del Proyecto Júpiter.}
\label{fig:arquitectura}
\end{figure}

\section{Modelo de Inteligencia Artificial e Hidrología}

El núcleo predictivo se basa en un modelo de \textit{Random Forest Classifier} entrenado con 21 variables meteorológicas e hidrológicas. Las variables de mayor peso determinante son:

\begin{enumerate}
    \item \textbf{Humedad del Suelo ($API_{72h}$):} Mide la cantidad de agua retenida en el terreno durante las últimas 72 horas.
    \item \textbf{Agua Acumulada en Superficie ($Q$):} Calculada mediante la ecuación del Número de Curva del Servicio de Conservación de Suelos (SCS-CN, USDA):
    \begin{equation}
        S = \frac{25400}{CN} - 254, \quad Q = \frac{(P - 0.2S)^2}{P + 0.8S}
    \end{equation}
    Donde $CN$ representa la impermeabilidad del suelo ($CN=88$ en cerros y $CN=90$ en zonas urbanas).
    \item \textbf{Pronósticos de Anticipación ($+1\text{h}, +3\text{h}, +6\text{h}$):} Vectores meteorológicos proyectados que permiten evaluar el riesgo horas antes de que la lluvia caiga físicamente.
\end{enumerate}

\section{Simulación Operativa: Evento del Viernes 31 de Julio}

Basado en el informe meteorológico para La Serena ($18.5\text{ mm}$ totales acumulados), se presenta el análisis detallado del comportamiento operativo durante la jornada.

\subsection{Análisis de los Pulsos de Precipitación}

La respuesta proyectada del modelo ante los principales pulsos de lluvia se resume en la Tabla~\ref{tab:pulsos}.

\begin{table}[htbp]
\centering
\caption{Resumen de predicciones del modelo para el evento del 31 de julio.}
\vspace{0.2cm}
\begin{tabular}{@{}ccllc@{}}
\toprule
\textbf{Hora Pulso} & \textbf{Precipitación} & \textbf{Hora Emisión Alerta} & \textbf{Llegada del Pico} & \textbf{Hora de Paso Seguro} \\ \midrule
11:00 hrs & 2.1 mm & 08:00 hrs (3h antes) & 11:00 hrs & 13:00 hrs \\
17:00 hrs & 5.4 mm & 14:00 hrs (3h antes) & 17:00 hrs & 19:30 hrs \\ \bottomrule
\end{tabular}
\label{tab:pulsos}
\end{table}

Como se observa en la Tabla~\ref{tab:pulsos}, el pulso de las 17:00 hrs ($5.4\text{ mm}$) resulta ser el más crítico del día. Esto se debe a que la lluvia previa de la mañana ($12.8\text{ mm}$ acumulados hasta las 14:00 hrs) habrá saturado el terreno. Por ende, la precipitación de las 17:00 hrs se convertirá casi en su totalidad en agua acumulada en superficie ($Q$). El modelo emitirá la alerta a las 14:00 hrs, otorgando una ventana de 3 horas para el resguardo de la población.

La Figura~\ref{fig:linea_tiempo} resume la secuencia temporal completa.

\begin{figure}[htbp]
\centering
\begin{tikzpicture}[xscale=1.9, yscale=0.8, >=Stealth]
    % Línea temporal
    \draw[->, line width=1.2pt] (0,0) -- (6.6,0) node[right] {Tiempo};
    
    % Hitos temporales bien alineados
    \foreach \x/\h/\p in {
        0.5/05:00/0.9 mm,
        1.5/08:00/0.5 mm,
        2.8/11:00/2.1 mm,
        4.0/14:00/4.9 mm,
        5.2/17:00/5.4 mm,
        6.2/20:00/0.3 mm
    } {
        \draw[line width=1.1pt] (\x,0.15) -- (\x,-0.15) node[below=2pt, font=\small] {\h};
        \node[font=\small\bfseries] at (\x, 0.6) {\p};
    }

    % Anotaciones explicativas superiores
    \node[font=\tiny, align=center] at (0.5, 1.3) {Lluvia inicial};
    \node[font=\tiny, align=center] at (1.5, 1.3) {Alerta 11:00};
    \node[font=\tiny, align=center] at (2.8, 1.3) {Primer pulso};
    \node[font=\tiny, align=center] at (4.0, 1.3) {Alerta Máxima};
    \node[font=\tiny, align=center] at (5.2, 1.3) {Pico del temporal};
    \node[font=\tiny, align=center] at (6.2, 1.3) {Inicio Calma};
\end{tikzpicture}
\caption{Línea de tiempo de procesamiento del modelo para el viernes 31 de julio.}
\label{fig:linea_tiempo}
\end{figure}

\section{Conclusiones y Sustento Técnico}

El Proyecto Júpiter demuestra que es posible transicionar de un esquema de respuesta reactivo a un modelo de prevención predictiva. Mediante el uso de métricas claras como la \textbf{Llegada del Pico} y la \textbf{Hora de Paso Seguro (Calma)}, las autoridades de emergencia (Bomberos y COGRID) pueden tomar decisiones informadas, garantizando la seguridad de la comunidad en La Serena.

\begin{thebibliography}{9}
\bibitem{usda1986} United States Department of Agriculture (USDA), \textit{Urban Hydrology for Small Watersheds}, Technical Release 55 (TR-55), 1986.
\bibitem{breiman2001} L. Breiman, \textit{Random Forests}, Machine Learning, vol. 45, no. 1, pp. 5--32, 2001.
\bibitem{openmeteo} Open-Meteo Weather API, \textit{Non-Real-Time Data Ingestion Protocol}, 2026.
\end{thebibliography}

\end{document}
